\section{Discussion}
\label{sec:discussion}

\paragraph{What sealing buys, and what it does not.}
Sealing converts a skill from a file anyone can copy into an artifact that only named providers can
open, and it keeps that file off every client. That closes the cheapest attack, copying the folder,
and moves the contest to the model's output channel. There the evidence in \cref{sec:eval} is
sobering: a skill that is in context can be extracted by a determined user unless the provider
adds defences, and even with defences some content leaks through transformation and paraphrase.
Sealing is therefore best understood as protection against \emph{redistribution}, plus a
measurable reduction in \emph{extraction}, not as secrecy of the skill's content.

\paragraph{Keep crown jewels out of context.}
The one design that provides secrecy against $\mathcal{A}_{\text{query}}$ is to keep the valuable
logic out of the model's context and sandbox altogether, behind an owner-hosted tool that returns
conclusions rather than rules (\cref{sec:casestudy}). The sealed skill then carries only when and
how to call the tool. This splits a skill into a \emph{protectable-from-copying} part (the
instructions) and a \emph{secret} part (the tool), and we expect commercial skills to converge on
this split.

\paragraph{Utility is the price of every output defence.}
Skills are useful precisely because the model acts on them, and acting on a rule often reveals it:
an answer that cites ``rule TXR-3: a difference above 0.50 but below 2\% is a soft mismatch'' is
both correct use and disclosure. Output formats that cite rule identifiers also hand the attacker
the vocabulary for targeted queries. Owners who seal a skill should write it so that its outputs
cite conclusions, not rule text, and providers should report the benign block rate of any filter
they deploy alongside its leak reduction.

\paragraph{Why no handshake.}
Dropping the onboarding handshake makes \sys{} as easy to adopt as encrypting to a published key,
and lets a provider support many owners without bilateral agreements. The costs are real: no
signed terms, no audit feed and no revocation beyond not renewing licences. These can be layered
on later without changing the package format, which is why we keep the base design minimal.

\paragraph{Artifact versus function.}
Everything \sys{} protect is an \emph{artifact}: the skill's text, code and data. Recent skill-stealing
attacks reconstruct a skill's \emph{function} from ordinary task results without ever asking for its
text~\citep{tsai2026daydreaming,hua2026behavioral}. No artifact-level protection, including an
owner-hosted tool, prevents that, because selling task results is the product. Sealing should be
judged by what it can deliver: the owner's verbatim work cannot be copied, redistributed or loaded
into a competing product, and a leaked reconstruction can be attributed.

\paragraph{Auditability.}
A sealed skill is opaque to the customer who runs it, and skills are an established injection and
supply-chain vector~\citep{schmotz2026skillinject,xu2026agentskills}. The provider is the one party
that can open every sealed skill, so it is the natural place to scan them before execution, and the
owner's signature gives provenance. Regulated customers may still need a named auditor, which the
format supports as one more recipient with its own published key.

\paragraph{Limitations.}
Our extraction study uses two 7--8B open models run locally with greedy decoding and a single
sample per prompt; frontier hosted models have stronger instruction-following and may be both
more resistant to naive requests and more capable of faithful paraphrase. The corpus is eight
synthetic skills of 290--360 words (eight to nine rules each) written for this study, and the 24 attacks are hand-written
rather than optimised. The semantic filter's threshold is calibrated on the same benign prompts on
which its block rate is reported, and benign rates rest on 32 tasks per configuration. Rule
disclosure is judged by a 7B model constrained by mechanical evidence checks; our manual check of
its decisions is small (85 decisions) and has not been repeated by an independent annotator. The tool-use case study is
illustrative (six attacks and two probes per cell), and its outcomes depend on how reliably each
model issues tool calls. Multi-turn filtering is evaluated post hoc: in a live deployment
a blocked turn changes the rest of the conversation. The design has not been implemented by any AI
provider; we are not aware of any provider that currently publishes a skill-sealing key or loads
encrypted skills, so our provider is a mock.
